\begin{ex}
Đơn vị đo góc thường dùng là gì? (Độ ghi 1, mét ghi 2).
\shortans{1}
\loigiai{
Đơn vị đo góc là độ ($^\circ$).
}
\end{ex}

\begin{ex}
Góc vuông có số đo bằng bao nhiêu độ?
\shortans{90}
\loigiai{
Góc vuông có số đo bằng $90^\circ$.
}
\end{ex}

\begin{ex}
Góc bẹt có số đo bằng bao nhiêu độ?
\shortans{180}
\loigiai{
Góc bẹt có số đo bằng $180^\circ$.
}
\end{ex}

\begin{ex}
Dụng cụ nào thường dùng để đo số đo của một góc? (Thước đo góc ghi 1, compa ghi 2).
\shortans{1}
\loigiai{
Thước đo góc (thước đo độ).
}
\end{ex}

\begin{ex}
Một nửa góc vuông có số đo bằng bao nhiêu độ?
\shortans{45}
\loigiai{
$90 : 2 = 45^\circ$.
}
\end{ex}

\begin{ex}
Góc bẹt có số đo gấp mấy lần góc vuông?
\shortans{2}
\loigiai{
$180 : 90 = 2$ lần.
}
\end{ex}

\begin{ex}
Góc tạo bởi hai kim đồng hồ lúc 3 giờ đúng có số đo bao nhiêu độ?
\shortans{90}
\loigiai{
Lúc 3 giờ đúng tạo thành góc vuông $90^\circ$.
}
\end{ex}

\begin{ex}
Góc tạo bởi hai kim đồng hồ lúc 6 giờ đúng có số đo bao nhiêu độ?
\shortans{180}
\loigiai{
Lúc 6 giờ đúng tạo thành góc bẹt $180^\circ$.
}
\end{ex}

\begin{ex}
Góc tạo bởi hai kim đồng hồ lúc 9 giờ đúng có số đo bao nhiêu độ?
\shortans{90}
\loigiai{
Lúc 9 giờ đúng tạo thành góc vuông $90^\circ$.
}
\end{ex}

\begin{ex}
Góc tạo bởi hai kim đồng hồ lúc 12 giờ đúng có số đo bao nhiêu độ?
\shortans{0}
\loigiai{
Hai kim trùng nhau tạo thành góc $0^\circ$.
}
\end{ex}

\begin{ex}
Mỗi góc nhọn của một chiếc ê-ke cân có số đo bằng bao nhiêu độ?
\shortans{45}
\loigiai{
Hai góc nhọn của ê-ke vuông cân bằng $45^\circ$.
}
\end{ex}

\begin{ex}
Trong tam giác đều, mỗi góc có số đo bằng bao nhiêu độ?
\shortans{60}
\loigiai{
Mỗi góc trong tam giác đều bằng $60^\circ$.
}
\end{ex}

\begin{ex}
Tổng số đo của hai góc vuông là bao nhiêu độ?
\shortans{180}
\loigiai{
$90 + 90 = 180^\circ$.
}
\end{ex}

\begin{ex}
Tổng số đo của 4 góc vuông trong một hình chữ nhật là bao nhiêu độ?
\shortans{360}
\loigiai{
$90 \times 4 = 360^\circ$.
}
\end{ex}

\begin{ex}
Góc có số đo lớn hơn $0^\circ$ và bé hơn $90^\circ$ gọi là góc gì? (Góc nhọn ghi 1, góc tù ghi 2).
\shortans{1}
\loigiai{
Góc nhọn.
}
\end{ex}

\begin{ex}
Góc có số đo lớn hơn $90^\circ$ và bé hơn $180^\circ$ gọi là góc gì? (Góc nhọn ghi 1, góc tù ghi 2).
\shortans{2}
\loigiai{
Góc tù.
}
\end{ex}

\begin{ex}
Một góc có số đo $60^\circ$, nếu tăng thêm $30^\circ$ thì trở thành góc gì? (Góc vuông ghi 1, góc bẹt ghi 2).
\shortans{1}
\loigiai{
$60 + 30 = 90^\circ$ là góc vuông.
}
\end{ex}

\begin{ex}
Một góc có số đo $120^\circ$, thêm bao nhiêu độ nữa thì bằng góc bẹt?
\shortans{60}
\loigiai{
$180 - 120 = 60^\circ$.
}
\end{ex}

\begin{ex}
Một góc vuông bị chia thành hai góc, trong đó một góc có số đo $35^\circ$. Góc còn lại có số đo bao nhiêu độ?
\shortans{55}
\loigiai{
$90 - 35 = 55^\circ$.
}
\end{ex}

\begin{ex}
Một góc bẹt bị chia thành hai góc, trong đó một góc có số đo $110^\circ$. Góc còn lại có số đo bao nhiêu độ?
\shortans{70}
\loigiai{
$180 - 110 = 70^\circ$.
}
\end{ex}

\begin{ex}
Kim phút quay được một phần tư vòng tròn quanh mặt đồng hồ thì tạo thành một góc bao nhiêu độ?
\shortans{90}
\loigiai{
Một phần tư vòng là góc $90^\circ$.
}
\end{ex}

\begin{ex}
Kim phút quay được nửa vòng tròn quanh mặt đồng hồ thì quét một góc bao nhiêu độ?
\shortans{180}
\loigiai{
Nửa vòng là góc bẹt $180^\circ$.
}
\end{ex}

\begin{ex}
Kim giờ quay từ 1 giờ đến 2 giờ thì quét một góc bao nhiêu độ?
\shortans{30}
\loigiai{
$360 : 12 = 30^\circ$.
}
\end{ex}

\begin{ex}
Lúc 2 giờ đúng, góc tạo bởi kim giờ và kim phút có số đo bao nhiêu độ?
\shortans{60}
\loigiai{
$30 \times 2 = 60^\circ$.
}
\end{ex}

\begin{ex}
Lúc 4 giờ đúng, góc tạo bởi kim giờ và kim phút có số đo bao nhiêu độ?
\shortans{120}
\loigiai{
$30 \times 4 = 120^\circ$.
}
\end{ex}
